The subject of K-Theory was historically concerned with studying vector bundles on spaces.
In Algbraic $K$-Theory, the primary objects of interest are projective modules over a ring, and vector bundles over
schemes.

\section{Projective Modules}

  Let $R$ be a (unital) associative ring. We mention the standard notion of projective modules  
  and illustrate its relevance to the subject of $K$-theory.

  Projective modules are closely related to free modules, and in fact.

\begin{definition}[Projective Module]
  We say $P$ is a projective (right\footnote{we fix a convention, either right or left}) $R$-module if it satisfies the following property:
    Let $M \to N$ be a surjection of (right) $R$-modules, then any map $\varphi : P \to N$ lifts to
  a map $\widetilde{\varphi} : P \to M$ such that the following diagram commutes:
  \[ 
   \begin{tikzcd}[row sep=large, column sep=large]
                     & P \arrow[dl, dashed, "\exists\widetilde\varphi"'] \arrow[d, "\varphi"] & \\
    M \arrow[r, two heads, "f"'] & N
    \end{tikzcd}
  \]
\end{definition}

It is a standard result that this definition is equivalent to the various following criteria:
\begin{proposition}
Let $R$ be as before and $P$ be an (right) $R$-module. Then the following are equivalent: 
\begin{itemize}
  \item [(i)] $P$ is projective. 
  \item [(ii)] For any surjection $M \adepiback N$ of (right) $R$-modules, the induced map $\Hom(P,M) \to \Hom(P,N)$
  is a surjection.
  \item [(iii)] Any surjection $\pi : M \adepiback P$ of (right) $R$-modules splits, i.e., there is $s : P \to M$ such that $\pi \circ s = 1_P$.
  \item [(iv)] There exists an $R$-module $Q$ such that $P \oplus Q$ is a free $R$-module.
\end{itemize}
\end{proposition}

%\begin{proof}
%  
%\end{proof}

\subsection{Basic Properties of Projective Modules}

When the ring $R$ is commutative, right $R$-modules can simply be called $R$-modules.

\begin{example}
All free modules are obviously projective. 
\end{example}

We claim that there exist non-free projective modules, and the failure of all projective modules
to be free shall be a major core of the theory henceforth.

\begin{example}\label{ch1:component_free}
Let $R$ and $S$ be non-trivial commutative rings. Then $R \oplus 0$ and $0 \oplus S$ are unsurprisingly projective $R \oplus S$-modules since $R \oplus S$ is a free
$R \oplus S$ module.
Yet, neither is free.
\end{example}

It is a question of interest to characterize when are projectives are free. 
\begin{proposition}
If $R$ is a field or a division ring, then all (right) projective modules over $R$ are free. 
\end{proposition}
%\begin{proof}

%\end{proof}

We also have the following result for commutative rings $R$: 

\begin{proposition}
All projective modules over a (commutative) local ring $R$ are free.
\end{proposition}
It is a standard exercise in algebra to show this is true for finitely generated projective modules.
However, the result holds more generally for any projective module over $R$ which we refer the reader to [Kaplansky].

\begin{proposition}\label{prop:projlocal}
If $\p$ is a prime ideal of a commutative ring $R$ and $P$ is a finitely generated projective $R$-module,
then the localization $P_p$ is free, that is, $P_p \cong (R_p)^n$ for some $n$.

Furthermore, a (finite) subcollection of the open sets $D(s)$ of $\Spec(R)$ may be chosen, say $\{D(s_i)\}$, to cover $R$
such that $P_{s_i} \cong (R[\frac{1}{s_i}])^{n_i}$ for each $i$.
\end{proposition}

\begin{proposition}
All projective modules over an Euclidean domain are free.
\end{proposition}


\section{Constructing Projective Modules}
To construct more interesting examples of projective modules, we have the following 
``recipes" inspired from the contruction of a vector bundle in the topological case.

\subsection{Open Patching}
Let $R$ be a commutative ring. Let $s_1, \ldots s_n \in R$ be such that 
$s_1R + \cdots + s_nR = R$. That is, $D(s_i) \cong \text{Spec}(R[\frac{1}{s_i}])$ cover $\text{Spec}(R)$.
Suppose we are given $g_{ij} \in GL_n(R[\frac{1}{s_is_j}])$ following the \emph{cocycle conditions:}

\begin{enumerate}
\item $g_{ii} = 1$ for all $i$.
  \item $g_{ij} g_{jk} = g_{ik}$ in $GL_n(R[\frac{1}{s_i s_j s_k}])$ for every $i,j,k$.
\end{enumerate}

Then 
\[
\  P = \left\{(x_1, \ldots, x_c) \in \prod\limits_{i = 1}^c \left(R\left[\frac{1}{s_i}\right]\right)^n : g_{ij}(x_j) = x_i \text{ in } R\left[\frac{1}{s_is_j}\right]^n \text{ for all } i, j\right\} 
\]
is a finitely generated projective $R$-module.

%\section{Idempotence Criteria for Projective Modules}
%Observe that for a unital ring $R$ and a projective (right) $R$-module $P$,
%suppose $Q$ is such that $P \oplus Q \cong R^n$.
%This describes an idempotent element $e \in M_n(R)$ such that $P = e(R^n)$ and.
%Any other choice of an idempotent $f$ such that $e(R^n) = f(R^n)$ leads to a conjugation

\section{Swan's Theorem}

The following result due to Swan \cite{Swan} shows why projective modules are the correct 
algebraic analogue to vector bundles in the topological setting. 

\begin{theorem}
  Let $\mathbb F$ be either $\mathbb R$ or $\mathbb C$.
  Let $X$ be a compact Hausorff space, and write $R$ for $C^0(X,\mathbb F)$. Consider the \emph{global sections functor} 
  $\Gamma : E \mapsto \Gamma(E,\mathbb F)$ from vector bundles over $X$ to projective modules over $R$.
  Then $\Gamma$ is an equivalence of categories.
\end{theorem}

The algebraic analogue of vector bundles are \emph{locally free sheaves over schemes},
which are an enrichment of the category of projective modules over unital rings.

